%\chapter*{Введение}
\chap{Введение}

Целью данной лабораторной работы является изучение влияния способа организации данных на скорость выполнения основных операций, освоение доступа к данным по индексу и по ключу, понимание принцип работы хеш-таблиц и причины возникновения коллизий, обучение работы с бинарной кучей и приоритетной очередью, а также реализация поиск Top-K элементов без полной сортировки данных.
\newline

В ходе выполнения лабораторной работы необходимо:
\begin{enumerate}
	\item сравнить стоимость различных способов доступа к данным;
	\item реализовать линейный поиск и простую хеш-таблицу с обработкой коллизий;
	\item исследовать влияние коэффициента заполнения;
	\item самостоятельно реализовать операции sift up и sift down для бинарной кучи;
	\item сравнить два способа поиска нескольких наибольших элементов.
\end{enumerate}

\chap{Выполнение}

\sect{Задание 1. Создание проекта}

\begin{figure}[!htbp]
	\centering
	\includegraphics[width=0.6\linewidth]{"Screenshot 2026-09-28 031606"}
	\caption{Структура проекта}
\end{figure}
\FloatBarrier

\sect{Задание 2. Доступ по индексу и линейный поиск}

\lstinputlisting[style=python, caption={}]{main.py}

\sect{Задание 3. Доступ по ключу}

\sect{Задание 4. Исследование функции hash()}

\sect{Задание 5. Простая хеш-таблица методом цепочек}

\lstinputlisting[style=python, caption={}]{hash_table.py}

\sect{Задание 6. Коэффициент заполнения и коллизии}

\sect{Задание 7. Представление бинарной кучи}

\sect{Задание 8. Реализация sift up и вставки}

\lstinputlisting[style=python, caption={}]{binary_heap.py}

\sect{Задание 9. Реализация sift down и извлечения корня}

\sect{Задание 10. Приоритетная очередь}

\sect{Задание 11. Top-K с полной сортировкой}

\sect{Задание 12. Top-K с помощью кучи}

\sect{Задание 13. Самостоятельная работа}

%\lstinputlisting[style=python, caption={}]{}

\chap{Вывод}

В ходе лабораторной работы:
объяснять, почему способ организации данных влияет на скорость выполнения операций;
различать доступ по известному индексу и последовательный поиск значения;
использовать словарь для доступа по ключу и объяснять его связь с хеш-таблицей;
объяснять назначение хеш-функции и причину возникновения коллизий;
реализовывать простую хеш-таблицу методом цепочек;
вычислять коэффициент заполнения хеш-таблицы и анализировать его влияние;
представлять бинарную кучу в массиве и вычислять индексы родителя и потомков;
реализовывать операции sift up и sift down;
использовать бинарную кучу как основу приоритетной очереди;
объяснять различие между полной сортировкой и поиском Top-K;
оценивать сложность основных операций и выбирать структуру данных исходя из требуемых операций.

\chap{Контрольные вопросы}

\begin{enumerate}
	\item \textit{Почему выбор структуры данных влияет на скорость программы?}\newline
	Скорость операций поиска, добавления и удаления во многом зависит от того, как именно организован доступ к элементам в структуре данных.
	\item \textit{Что означает асимптотическая сложность O(1)?}\newline
	Оценка O(1) означает, что количество основных шагов алгоритма не растёт пропорционально размеру набора данных.
	\item \textit{Почему доступ к элементу списка по известному индексу считается операцией O(1)?}\newline
	В списке Python элементы располагаются в определенном порядке и имеют индексы. Если индекс известен, доступ к элементу выполняется напрямую.
	\item \textit{Как работает линейный поиск и какова его сложность в худшем случае?}\newline
	Простейший линейный поиск последовательно сравнивает элементы с искомым значением. В худшем случае алгоритму приходится просмотреть весь список, следовательно сложность O(n).
	\item \textit{Чем доступ по индексу отличается от доступа по ключу?}\newline
	Поиск по ключу реализован на основе хеш-таблицы в среднем выполняется значительно быстрее, чем последовательный. Но из-за дополнительной памяти и вычисления хеша ключа его внутреннее устройство сложнее.
	\item \textit{Для каких задач удобно использовать словарь dict?}\newline
	Если имеется у набора объектов имеются некие индивидуальные ключи.
	\item \textit{Что такое хеш-функция?}\newline
	Функция, реализующая алгорим хэширования -- способа преобразования ключа в числовое значение, которое используется для определения позиции элемента в хеш-таблице.
	\item \textit{Для чего хеш ключа преобразуется в индекс корзины?}\newline
	\item \textit{Что называется коллизией в хеш-таблице?}\newline
	Коллизией называется ситуация, когда разные ключи после применения хеш-функции попадают в одну и ту же позицию таблицы.
	\item \textit{Почему коллизии нельзя считать ошибкой реализации?}\newline
	Коллизии являются нормальным свойством хеш-таблиц, и полностью исключить их для произвольного множества ключей обычно невозможно.
	\item \textit{В чём заключается метод цепочек?}\newline
	Каждая ячейка таблицы хранит не один элемент, а список пар ключ--значение. Для поиска элемента сначала вычисляется индекс корзины, после чего просматриваются только элементы соответствующей цепочки.
	\item \textit{Что показывает коэффициент заполнения хеш-таблицы?}\newline
	Отношение числа сохраненных элементов к количеству корзин. 
	\item \textit{Как увеличение коэффициента заполнения может повлиять на поиск?}\newline
	По мере увеличения коэффициента заполнения вероятность длинных цепочек возрастает, и поиск начинает требовать больше сравнений.
	\item \textit{Что такое бинарная куча?}\newline
	\item \textit{Какое свойство выполняется в min-heap?}\newline
	Структура данных, представляющая собой почти полное бинарное дерево и обычно хранящаяся в обычном массиве
	\item \textit{Почему минимальный элемент min-heap всегда находится в корне?}\newline
	Так как значение каждого родителя не больше значений его потомков.
	\item \textit{Как вычислить индекс родителя элемента кучи?}\newline
	$p = (i - 1) // 2$.
	\item \textit{Как вычислить индексы левого и правого потомков?}\newline
	$l = 2 * i + 1; r = 2 * i + 2$.
	\item \textit{Для чего используется операция sift up?}\newline
	Для восстановления свойства кучи после добавления нового элемента (т.к. он сначала помещается в конец массива) его последовательно сравнивают с родителем и при необходимости меняют их местами.
	\item \textit{Для чего используется операция sift down?}\newline
	Для восстановления свойства кучи после извлечения корневого элемента (т.к. на его место переносится последний элемент массива) его замену последовательно сравнивают с потомками и при необходимости опускают вниз.
	\item \textit{Какова сложность вставки элемента в бинарную кучу?}\newline
	Cложность составляет $O(\log_N)$.
	\item \textit{Какова сложность извлечения корня бинарной кучи?}\newline
	Cложность составляет $O(\log_N)$.
	\item \textit{Что такое приоритетная очередь?}\newline
	Структура данных, где у каждого добавленного элемента есть свой уровень приоритета, от которого зависит порядок обработки.
	\item \textit{Чем приоритетная очередь отличается от обычной очереди FIFO?}\newline
	Обычная очередь обрабатывает элементы в порядке поступления. В приоритетной очереди следующим извлекается элемент с наиболее высоким приоритетом, независимо от того, когда он был добавлен.
	\item \textit{Как модуль heapq интерпретирует вершину кучи?}\newline
	В стандартном модуле heapq первым извлекается элемент с наименьшим значением. Если приоритет задается числом, необходимо заранее определить, означает ли меньшее число более высокий приоритет.
	\item \textit{Что означает задача Top-K?}\newline
	Во многих задачах требуется найти не полный отсортированный набор, а только K лучших элементов.
	\item \textit{Почему для поиска Top-K не всегда требуется сортировать весь набор данных?}\newline
	Куча гарантирует только положение корневого элемента и выполнение свойства кучи.
	\item \textit{Как работает алгоритм Top-K с min-heap размером K?}\newline
	Если K значительно меньше n, можно поддерживать min-heap размером K. После обработки всех данных в куче останутся K наибольших значений.
	\item \textit{Какова сложность Top-K через кучу размера K?}\newline
	Cложность составляет $O(n\log_K)$, т.к. для каждого нового элемента выполняется не более одной операции над кучей размера K.
	\item \textit{Почему элементы внутри бинарной кучи не обязаны быть полностью отсортированы?}\newline
	Бинарная куча по определению имеет определенный порядок хранения элементов.
	\item \textit{В каком случае полная сортировка может быть проще и оправданнее, чем поддержание кучи?}\newline
	При малом количестве элементов поддержание кучи имеет большую сложность, чем полная сортировка.
	\item \textit{Какие компромиссы между скоростью и памятью возникают при использовании хеш-таблиц?}\newline
	Быстрый доступ достигается за счет дополнительной памяти и вычисления хеша ключа.
\end{enumerate}
